\chapter{有限 inner 求解、outer 同步与步长}
\label{chap:outer}

\section{以振幅定义误差}
固定 teacher，定义
\[
 d_{\mathrm{proj}}(p,q)=\osc_a\log\frac{p(a)}{q(a)},\qquad
 \osc_a u=\max_a u(a)-\min_a u(a).
\]
这是严格正概率 simplex 上的 Hilbert projective 距离。
加动作无关常数不改变振幅；三角不等式来自
$\osc(u+v)\le\osc u+\osc v$。
若 $\osc\log(p/q)=0$，则 $p=cq$，归一化要求 $c=1$。
采用这个定义后，inner 误差项前面没有额外系数 $2$。

\begin{theorem}[exact outer 与所有步长区间]\label{thm:projective}
固定 $T>0$，令 $p_{k+1}=G_{\eta_k}(p_k)$。则
\[
 \Delta_{k+1}=|1-\eta_k|\Delta_k,\qquad
 \Delta_k=d_{\mathrm{proj}}(p_k,T).
\]
对常数步长，$0<\eta<1$ 同向收敛，$\eta=1$ 一步投影，
$1<\eta<2$ 差分交替收敛，$\eta=2$ 振幅不变。
$\eta>2$ 且 $p_0\ne T$ 时振幅发散。
\end{theorem}
\begin{proof}
由 Gibbs 公式，
\[
 \log\frac{p_{k+1}(a)}{T(a)}
 =(1-\eta_k)\log\frac{p_k(a)}{T(a)}-\log Z_k.
\]
对任意动作对相减，$\log Z_k$ 消失。对实数 $c$，
$\osc(cu)=|c|\osc u$，代入即得恒等式。
$\Delta_k\to0$ 推出 $p_k\to T$：令 $u_a=p_k(a)/T(a)$，
归一化 $\sum_a T(a)u_a=1$ 要求 $\min u\le1\le\max u$，
故 $|\log u_a|\le\Delta_k$，每个坐标收敛。
\end{proof}

\section{提前同步后重新 rollout}
第 $k$ 轮 exact target 为 $q_{k+1}=G_{\eta_k}(p_k)$，
实际 student 为 $p_{k+1}$，误差定义为
\[
 \varepsilon_k=d_{\mathrm{proj}}(p_{k+1},q_{k+1}).
\]
下一轮保存的是实际 $p_{k+1}$，信号为 $\log T-\log p_{k+1}$，
目标为 $G_{\eta_{k+1}}(p_{k+1})$，不会继续优化旧 target。

\begin{theorem}[非精确 outer 的误差传播]\label{thm:inexact}
若 $0<\eta_k\le1$，则
\[
 \Delta_K\le P_{0,K}\Delta_0+
 \sum_{j=0}^{K-1}P_{j+1,K}\varepsilon_j,\qquad
 P_{i,K}=\prod_{\ell=i}^{K-1}(1-\eta_\ell).
\]
常数 $\eta$ 且 $\varepsilon_k\le\bar\varepsilon$ 时
$\limsup_K\Delta_K\le\bar\varepsilon/\eta$。
常数 $\eta>0$ 且 $\varepsilon_k\to0$ 时 $\Delta_K\to0$。
\end{theorem}
\begin{proof}
由三角不等式，
\[
 \Delta_{k+1}\le d_{\mathrm{proj}}(p_{k+1},q_{k+1})
             +d_{\mathrm{proj}}(q_{k+1},T)
 =\varepsilon_k+(1-\eta_k)\Delta_k.
\]
代入前一轮，再逐轮展开，得到乘积加权式。
常数步长下几何级数和为
$\sum_{j=0}^{K-1}(1-\eta)^j=[1-(1-\eta)^K]/\eta$。
证明误差趋零时，任取 $a>0$，选择 $N$ 使 $j\ge N$ 时
$\varepsilon_j\le a\eta$。有限个 $j<N$ 的项随 $K$ 衰减为零，
剩余尾和至多为 $a$；令 $a\downarrow0$ 完成证明。
\end{proof}
第一轮没收敛的影响是 $P_{1,K}\varepsilon_0$，会被后续正确更新逐渐遗忘。
新 rollout 噪声与未访问 prefix 的误差仍进入新的 $\varepsilon_k$。
该递推需要逐 prefix 的误差控制，不由单个 batch 平均 loss 自动推出。

\section{固定 anchor 的参数局部模型}
令 logits Jacobian 为 $J_k(s)$，categorical Fisher 为
$C_k(s)=\diag(p_k)-p_kp_k^\top$，固定训练测度下
\[
 A_k=\E_s[J_k(s)^\top C_k(s)J_k(s)],\qquad
 Q_k(\delta)=\frac1{2\eta_k}\delta^\top A_k\delta-b_k^\top\delta.
\]
这是把 reward 项一阶线性化、KL 项二阶展开后的模型，
不是一般真实目标完整 Hessian；真实 reward Hessian 会增加 $O(\|\delta\|^2)$ 项。
在可辨识商空间 $A_k$ 的谱位于 $[m_k,L_k]$，最优位移为
$\delta_k^\star=\eta_k A_k^{-1}b_k$。

\begin{theorem}[局部 GD 与 Natural-GD 的 inner 精度]\label{thm:inner}
普通 GD 使用 $0<\gamma<2\eta_k/L_k$，则
\[
 \|e_M\|_{A_k}\le q_k^M\|e_0\|_{A_k},\qquad
 q_k=\max_{\lambda\in[m_k,L_k]}|1-\gamma\lambda/\eta_k|.
\]
最优常数步长为 $2\eta_k/(m_k+L_k)$，对应
$q_k=(L_k-m_k)/(L_k+m_k)$。
预条件步 $\delta_{m+1}=\delta_m-\gamma A_k^{-1}\nabla Q_k(\delta_m)$
满足 $e_{m+1}=(1-\gamma/\eta_k)e_m$，$\gamma=\eta_k$ 时一步解模型。
\end{theorem}
\begin{proof}
最优性给出 $A_k\delta_k^\star=\eta_kb_k$，故
\[
 e_{m+1}=(I-\gamma A_k/\eta_k)e_m.
\]
用正交矩阵对角化 $A_k=U\Lambda U^\top$，
\[
 \|e_{m+1}\|_{A_k}^2
 =\sum_i\lambda_i(1-\gamma\lambda_i/\eta_k)^2(U^\top e_m)_i^2
 \le q_k^2\|e_m\|_{A_k}^2.
\]
递推得到第一式。端点绝对乘子相等时最坏谱半径最小，
解 $1-\gamma m_k/\eta_k=-(1-\gamma L_k/\eta_k)$ 得最优步长。
预条件时直接消去 $A_k$，得到标量乘子。
\end{proof}
病态 Fisher 增加 inner 数量。实际 Natural-GD 用 damping 或采样 Fisher，
需分析 $P^{1/2}A_kP^{1/2}$ 的谱；伪逆只在 $b_k$ 位于像空间且去掉 gauge 后有效。

若在局部邻域有
$d_{\mathrm{proj}}(\pi_\theta,\pi_{\theta'})\le C_k\|\theta-\theta'\|_{A_k}$，
才可将参数 inner 误差转为
\[
 \varepsilon_k\le C_kq_k^{M_k}\|\delta_k^\star\|_{A_k}
              +C_{\mathrm{nl}}\|\delta_k^\star\|^2
              +\varepsilon_k^{\mathrm{sampling}}
              +\varepsilon_k^{\mathrm{representation}}.
\]
增大 $M_k$ 只能减少第一项。$\eta=1$ 时 target 始终是 teacher，
有限 inner 多轮收敛来自求解器收缩，重新 rollout 负责修正覆盖。

\section{Hard、delayed 与逐步同步}
Hard sync 在 batch 边界复制 student，并冻结 $r_k,T,p_k$ 直到 inner 结束，
这是主算法。停止条件应控制到本轮 target 的误差。
若要求 $\varepsilon_k\le c\eta_k^2$，常数步长稳态偏差为 $O(\eta_k)$；
若只要求 $\varepsilon_k\le c\eta_k$，稳态上界为 $O(1)$。
因此更小步长不自动提高精度。

Delayed sync 保持旧 anchor $\bar p$ 跨多批次。
固定 teacher 时有精确 target 比较
\[
 d_{\mathrm{proj}}(G_\eta(p),G_\eta(\bar p))
 =|1-\eta|\,d_{\mathrm{proj}}(p,\bar p).
\]
这量只度量 target drift；rollout distribution mismatch 还需另算。
若每个 optimizer step 都同步，旧 KL 在初始化点的一阶梯度反复为零。
将这些更新解释成固定子问题的 inner solve 会破坏定理假设；
若每步都解新的 exact 微小 mirror step，则是另一个有更多 outer 的算法。

\section{KL budget 与可变步长}
令 $r=\log(T/b)$，$A(\eta)=\log\E_b e^{\eta r}$，
则 $A'=\E_{G_\eta(b)}r$、$A''=\Var_{G_\eta(b)}r$。
\begin{theorem}[镜像路径 KL 的单调性]\label{thm:budget}
\[
 D(\eta)=\KL(G_\eta(b)\Vert b)=\eta A'(\eta)-A(\eta),\qquad
 D'(\eta)=\eta A''(\eta)\ge0.
\]
若 $b\ne T$，有限共同 support 下 $D$ 在 $\eta>0$ 严格递增。
\end{theorem}
\begin{proof}
归一化权重对 $\eta$ 求导为
$\partial_\eta p_\eta(a)=p_\eta(a)[r(a)-\E_{p_\eta}r]$，
所以 $A''=\E r^2-(\E r)^2$。代入 KL 表达式求导，$A'$ 相消。
$b\ne T$ 时 $r$ 非常数，正权重方差严格为正。
\end{proof}
预算 $\delta\in(0,\KL(T\Vert b))$ 在 $(0,1)$ 唯一可解；
预算大于该端点时应取 $\eta=1$，不能声称区间内总有解。
对固定 prefix，Taylor 展开
\[
 D(\eta)=\frac{\eta^2}{2}\Var_b(r)+O(\eta^3),\qquad
 \eta\simeq\sqrt{\frac{2\delta}{\widehat{\Var}_b(r)+\epsilon_{\mathrm{var}}}}.
\]
全局共享 $\eta$ 可用 $\nu$ 平均 KL 或高分位预算，但二者保障不同。
对 sampled actions，需要置信余量；单个经验方差不能排除未采样的尖锐动作。

\begin{corollary}[递减步长与 vanishing relative error]
若 $0<\eta_k<1$、$\sum_k\eta_k=\infty$ 且
$\varepsilon_k/\eta_k\to0$，则 $\Delta_k\to0$。
\end{corollary}
\begin{proof}
乘积满足 $P_{N,K}\le e^{-\sum_{k=N}^{K-1}\eta_k}\to0$。
任给 $a>0$，选 $N$ 使 $\varepsilon_k\le a\eta_k$（$k\ge N$）。
递推为 $\Delta_{k+1}-a\le(1-\eta_k)(\Delta_k-a)$，
所以 $\limsup\Delta_k\le a$。令 $a\downarrow0$。
\end{proof}
仅要求 $\varepsilon_k\to0$ 对递减 $\eta_k$ 不够；
如 $\varepsilon_k=c\eta_k$ 会留下半径 $c$。
固定小步长达到 $\epsilon$ 的 outer 数量约为
$\eta^{-1}\log(\Delta_0/\epsilon)$，应连同 rollout 成本和 inner 残差共同比较。

\section{线性步长调度的精确界}
固定 $\eta$ 不是 MirrorOPD 的必要条件。一个常用的实现是先用较小的镜像步，
再在线性调度中逐渐增大到目标值。令
\[
 \eta_k=\min\{\eta_{\max},\eta_{\min}+k\delta\},
 \qquad 0<\eta_{\min}\le\eta_{\max}\le1,
\]
其中 $\delta\ge0$。这里的调度只改变第 $k$ 个 outer 子问题的目标，
不改变该 outer 内固定 anchor 的定义。

\begin{theorem}[线性调度的乘积收缩与复杂度]
\label{thm:linear-schedule}
在精确 hard-sync、共同 support 且 $0<\eta_k\le1$ 的条件下，
\[
 d_{\mathrm{proj}}(p_K,T)
 =\Big(\prod_{k=0}^{K-1}(1-\eta_k)\Big)d_{\mathrm{proj}}(p_0,T)
 \le \exp\!\Big(-\sum_{k=0}^{K-1}\eta_k\Big)d_{\mathrm{proj}}(p_0,T).
\]
若在线性段内 $\eta_k=\eta_{\min}+k\delta$ 且
$\eta_{\min}+ (K-1)\delta\le1$，则
\[
 \sum_{k=0}^{K-1}\eta_k
 =K\eta_{\min}+\frac{\delta K(K-1)}2,
\]
从而达到 $d_{\mathrm{proj}}(p_K,T)\le\varepsilon$ 的充分条件是
\[
 K\eta_{\min}+\frac{\delta K(K-1)}2
 \ge\log\frac{d_{\mathrm{proj}}(p_0,T)}\varepsilon.
\]
在线性段结束后固定 $\eta_{\max}$，额外需要至多
\[
 \left\lceil \frac{\log(d_{\mathrm{proj}}(p_{K_0},T)/\varepsilon)}
 {-\log(1-\eta_{\max})}\right\rceil
\]
个 outer，其中 $K_0$ 是进入平台的时刻。
\end{theorem}

\begin{proof}
由定理~\ref{thm:projective}，每一个 exact outer 都满足
\[
 \Delta_{k+1}=(1-\eta_k)\Delta_k,
 \qquad \Delta_k=d_{\mathrm{proj}}(p_k,T).
\]
逐步相乘得到第一式。对 $x\in[0,1)$，不等式
$1-x\le e^{-x}$ 来自 $\log(1-x)\le-x$；对每个 $\eta_k$ 应用它，得到指数界。

在线性段内直接求等差数列的和：
\[
 \sum_{k=0}^{K-1}(\eta_{\min}+k\delta)
 =K\eta_{\min}+\delta\sum_{k=0}^{K-1}k
 =K\eta_{\min}+\frac{\delta K(K-1)}2.
\]
将其代入指数界并令右侧不超过 $\varepsilon$，得到充分条件。
平台阶段每一步乘以 $1-\eta_{\max}$，所以
\[
 \Delta_{K_0+j}\le(1-\eta_{\max})^j\Delta_{K_0}.
\]
取对数并向上取整即得最后一个上界。
\end{proof}

该定理说明线性调度的收益来自累计镜像强度，而不是来自改变 exact target 的收敛点。
在有限 inner 情形，若第 $k$ 轮误差为 $\varepsilon_k$，只需把定理
~\ref{thm:inexact} 的权重替换为
\[
 P_{j+1,K}=\prod_{\ell=j+1}^{K-1}(1-\eta_\ell),
 \qquad
 \Delta_K\le P_{0,K}\Delta_0+\sum_{j=0}^{K-1}P_{j+1,K}\varepsilon_j.
\]
因此增大 $\eta_k$ 同时带来两种相反效应：乘积收缩更快，但若 inner 预算不随
目标位移增加，$\varepsilon_k$ 可能变大。一个可验证的调度规则是令停止阈值满足
\[
 \varepsilon_k\le c\,\eta_k^{1+\rho},\qquad \rho>0,
\]
此时误差项相对于当轮镜像移动是渐近可忽略的；若只保持
$\varepsilon_k\le c\eta_k$，则只能保证一个与 $c$ 同阶的残差半径。
